In the first-order Drift-Kinetic equation, the momentum space
operator is
\begin{equation}
\nabla_{\mathbf{p}}\cdot\mathbf{S}_{\mathbf{p}}\left(  f_{1}\right)
\label{momentumdke}%
\end{equation}
where the fluxes are expressed as (\ref{momflux2}).
\begin{equation}
\mathbf{S}_{\mathbf{p}}\left(  f_{1}\right)  =-\mathbb{{D}_{\mathbf{p}}%
\cdot\nabla_{\mathbf{p}}f_{1}+\mathbf{F}_{\mathbf{p}}f_{1}}%
\end{equation}
with the components (\ref{momflux3})
\begin{align}
S_{p}  & =-D_{pp}\frac{\partial f_{1}}{\partial p}+D_{p\xi}\frac{\sqrt
{1-\xi^{2}}}{p}\frac{\partial f_{1}}{\partial\xi}+F_{p}f_{1}\\
S_{\xi}  & =-D_{\xi p}\frac{\partial f_{1}}{\partial p}+D_{\xi\xi}\frac
{\sqrt{1-\xi^{2}}}{p}\frac{\partial f_{1}}{\partial\xi}+F_{\xi}f_{1}%
\end{align}

The first-order distribution function $f_{1}$ is decomposed into
\begin{equation}
f_{1}=\widetilde{f}+g
\end{equation}
where $\widetilde{f}$ is a source term due to spacial gradients
\begin{equation}
\widetilde{f}=\frac{v_{\Vert}}{\Omega}I\left(  \psi\right)  \frac{\partial
f_{0}}{\partial\psi}%
\end{equation}
and $g$ is the response of the plasma.

All the operators in the conservative form above are readily linear, so that
we can decompose (\ref{momentumdke}) as
\begin{equation}
\nabla_{\mathbf{p}}\cdot\mathbf{S}_{\mathbf{p}}\left(  f_{1}\right)
=\nabla_{\mathbf{p}}\cdot\mathbf{S}_{\mathbf{p}}\left(  \widetilde{f}\right)
+\nabla_{\mathbf{p}}\cdot\mathbf{S}_{\mathbf{p}}\left(  g\right)
\end{equation}

According to (\ref{kinbav2}), we can express the bounce-averaged operator as
\begin{align}
\left\{  J_{\mathbf{p}}\nabla_{\mathbf{p}}\cdot\mathbf{S}_{\mathbf{p}}\left(
\widetilde{f}\right)  \right\}   & =\frac{\partial}{\partial p}\left(
p^{2}\widetilde{S}_{p}^{\left(  0\right)  }\right)  -\frac{p}{\lambda}%
\frac{\partial}{\partial\xi_{0}}\left(  \sqrt{1-\xi_{0}^{2}}\lambda
\widetilde{S}_{\xi_{0}}^{\left(  0\right)  }\right) \\
\left\{  J_{\mathbf{p}}\nabla_{\mathbf{p}}\cdot\mathbf{S}_{\mathbf{p}}\left(
g\right)  \right\}   & =\frac{\partial}{\partial p}\left(  p^{2}S_{p}^{\left(
0\right)  }\right)  -\frac{p}{\lambda}\frac{\partial}{\partial\xi_{0}}\left(
\sqrt{1-\xi_{0}^{2}}\lambda S_{\xi_{0}}^{\left(  0\right)  }\right)
\label{kinbav3}%
\end{align}
where we need to evaluate the bounce-averaged fluxes (\ref{sp0})
\begin{align}
\widetilde{S}_{p}^{\left(  0\right)  }  & =\left\{  S_{p}\left(  \widetilde
{f}\right)  \right\} \\
\widetilde{S}_{\xi_{0}}^{\left(  0\right)  }  & =\sigma\left\{  \frac
{\sigma\xi}{\sqrt{\Psi}\xi_{0}}S_{\xi}\left(  \widetilde{f}\right)  \right\}
\end{align}
and
\begin{align}
S_{p}^{\left(  0\right)  }  & =\left\{  S_{p}\left(  g\right)  \right\} \\
S_{\xi_{0}}^{\left(  0\right)  }  & =\sigma\left\{  \frac{\sigma\xi}%
{\sqrt{\Psi}\xi_{0}}S_{\xi}\left(  g\right)  \right\}
\end{align}

Because $g$ is constant along a field line, we have $g\left(  p,\xi\right)
=g^{(0)}\left(  p,\xi_{0}\right)  $which is independent of $\theta$ and
$\sigma$. Therefore, the fluxes for $g$ have exactly the same expression as
for $f_{0}$ in the zero-order equation described in section . This is why they
have the same notation in (\ref{kinbav3}), while the fluxes associated with
$\widetilde{f}$ are noted $\widetilde{S}$. .

Indeed, $\widetilde{f}$ has an explicit dependence upon $\theta$, which can be
isolated as follows:
\begin{equation}
\widetilde{f}(p,\xi,\psi,\theta)=\frac{\xi(\psi,\theta,\xi_{0})}{\Psi
(\psi,\theta)\xi_{0}}\widetilde{f}^{(0)}(p,\xi_{0},\psi)\label{ftilde}%
\end{equation}
with
\begin{equation}
\widetilde{f}^{(0)}(p,\xi_{0},\psi)=\frac{p\xi_{0}I\left(  \psi\right)
}{q_{e}B_{\min}\left(  \psi\right)  }\frac{\partial f_{0}^{(0)}(p,\xi_{0}%
,\psi)}{\partial\psi}%
\end{equation}

We can note that $\widetilde{f}^{(0)}$ is antisymmetric in the trapped region,
since $f_{0}^{(0)}$ is symmetric and $\xi_{0}$ is, of course, anti-symmetric.
As a result, only $\sigma\widetilde{f}^{(0)}$ can be taken out of the bounce
averaging operator. Taking the bounce-average of each term, we find
\begin{align}
\left\{  -D_{pp}\frac{\partial\widetilde{f}}{\partial p}\right\}   &
=-\sigma\left\{  \sigma\frac{\xi}{\Psi\xi_{0}}D_{pp}\right\}  \frac
{\partial\widetilde{f}^{(0)}}{\partial p}\\
\left\{  D_{p\xi}\frac{\sqrt{1-\xi^{2}}}{p}\frac{\partial\widetilde{f}%
}{\partial\xi}\right\}   & =\frac{\sqrt{1-\xi_{0}^{2}}}{p}\left\{  \frac
{\xi^{2}}{\Psi^{3/2}\xi_{0}^{2}}D_{p\xi}\right\}  \frac{\partial\widetilde
{f}^{(0)}}{\partial\xi_{0}}+\frac{\sqrt{1-\xi_{0}^{2}}}{p}\sigma\left\{
\sigma\frac{\Psi-1}{\Psi^{3/2}\xi_{0}^{3}}D_{p\xi}\right\}  \widetilde
{f}^{(0)}\\
\left\{  F_{p}\widetilde{f}\right\}   & =\sigma\left\{  \sigma\frac{\xi}%
{\Psi\xi_{0}}F_{p}\right\}  \widetilde{f}^{(0)}\\
-\sigma\left\{  \frac{\sigma\xi}{\sqrt{\Psi}\xi_{0}}D_{\xi p}\frac
{\partial\widetilde{f}}{\partial p}\right\}   & =-\left\{  \frac{\xi^{2}}%
{\Psi^{3/2}\xi_{0}^{2}}D_{\xi p}\right\}  \frac{\partial\widetilde{f}^{(0)}%
}{\partial p}\\
\sigma\left\{  \frac{\sigma\xi}{\sqrt{\Psi}\xi_{0}}D_{\xi\xi}\frac{\sqrt
{1-\xi^{2}}}{p}\frac{\partial\widetilde{f}}{\partial\xi}\right\}   &
=\frac{\sqrt{1-\xi_{0}^{2}}}{p}\sigma\left\{  \frac{\sigma\xi^{3}}{\Psi^{2}%
\xi_{0}^{3}}D_{\xi\xi}\right\}  \frac{\partial\widetilde{f}^{(0)}}{\partial
\xi_{0}}+\frac{\sqrt{1-\xi_{0}^{2}}}{p}\left\{  \frac{\xi\left(
\Psi-1\right)  }{\Psi^{2}\xi_{0}^{4}}D_{\xi\xi}\right\}  \widetilde{f}^{(0)}\\
\sigma\left\{  \frac{\sigma\xi}{\sqrt{\Psi}\xi_{0}}F_{\xi}\widetilde
{f}\right\}   & =\left\{  \frac{\xi^{2}}{\Psi^{3/2}\xi_{0}^{2}}F_{\xi
}\right\}  \widetilde{f}^{(0)}%
\end{align}
where the following relation is used
\begin{align}
\frac{\partial}{\partial\xi_{0}}\frac{\xi}{\xi_{0}}  & =\sigma\frac{\partial
}{\partial\left|  \xi_{0}\right|  }\frac{\sqrt{1-\Psi\left(  1-\left|  \xi
_{0}\right|  ^{2}\right)  }}{\left|  \xi_{0}\right|  }\\
& =\sigma\frac{\Psi\xi_{0}^{2}-\xi^{2}}{\left|  \xi\right|  \xi_{0}^{2}}\\
& =\sigma\frac{\Psi-1}{\left|  \xi\right|  \xi_{0}^{2}}%
\end{align}

We can therefore rewrite
\[
\widetilde{\mathbf{S}}_{\mathbf{p}}^{(0)}\left(  \widetilde{f}^{(0)}\right)
=-\widetilde{\mathbb{{D}}}_{\mathbf{p}}^{(0)}\cdot\nabla_{p,\xi_{0}}%
\widetilde{f}^{(0)}+\widetilde{\mathbf{F}}_{\mathbf{p}}^{(0)}\widetilde
{f}^{(0)}%
\]
where the bounce averaged flux is decomposed into
\begin{equation}
\mathbf{S}_{\mathbf{p}}^{(0)}=\left(
\begin{array}
[c]{l}%
S_{p}^{(0)}\\
S_{\xi_{0}}^{(0)}%
\end{array}
\right)
\end{equation}
with
\begin{align}
\widetilde{S}_{p}^{(0)}  & =-\widetilde{D}_{pp}^{(0)}\frac{\partial
\widetilde{f}^{(0)}}{\partial p}+\frac{\sqrt{1-\xi_{0}^{2}}}{p}\widetilde
{D}_{p\xi}^{(0)}\frac{\partial\widetilde{f}^{(0)}}{\partial\xi_{0}}%
+\widetilde{F}_{p}^{(0)}\widetilde{f}^{(0)}\\
\widetilde{S}_{\xi_{0}}^{(0)}  & =-\widetilde{D}_{\xi p}^{(0)}\frac
{\partial\widetilde{f}^{(0)}}{\partial p}+\frac{\sqrt{1-\xi_{0}^{2}}}%
{p}\widetilde{D}_{\xi\xi}^{(0)}\frac{\partial\widetilde{f}^{(0)}}{\partial
\xi_{0}}+\widetilde{F}_{\xi}^{(0)}\widetilde{f}^{(0)}%
\end{align}
by defining the diffusion components
\begin{align}
\widetilde{D}_{pp}^{(0)}  & =\sigma\left\{  \sigma\frac{\xi}{\Psi\xi_{0}%
}D_{pp}\right\}  \label{diffbavdke}\\
\widetilde{D}_{p\xi}^{(0)}  & =\left\{  \frac{\xi^{2}}{\Psi^{3/2}\xi_{0}^{2}%
}D_{p\xi}\right\}  \\
\widetilde{D}_{\xi p}^{(0)}  & =\left\{  \frac{\xi^{2}}{\Psi^{3/2}\xi_{0}^{2}%
}D_{\xi p}\right\}  \\
\widetilde{D}_{\xi\xi}^{(0)}  & =\sigma\left\{  \frac{\sigma\xi^{3}}{\Psi
^{2}\xi_{0}^{3}}D_{\xi\xi}\right\}
\end{align}
and the convection components
\begin{align}
\widetilde{F}_{p}^{(0)}  & =\sigma\left\{  \sigma\frac{\xi}{\Psi\xi_{0}}%
F_{p}\right\}  +\frac{\sqrt{1-\xi_{0}^{2}}}{p\xi_{0}^{3}}\left\{
\frac{\left(  \Psi-1\right)  }{\Psi^{3/2}}D_{p\xi}\right\}  \\
\widetilde{F}_{\xi}^{(0)}  & =\left\{  \frac{\xi^{2}}{\Psi^{3/2}\xi_{0}^{2}%
}F_{\xi}\right\}  +\frac{\sqrt{1-\xi_{0}^{2}}}{p\xi_{0}^{3}}\sigma\left\{
\frac{\sigma\xi\left(  \Psi-1\right)  }{\xi_{0}\Psi^{2}}D_{\xi\xi}\right\}
\label{convbavdke}%
\end{align}
where we use that $\sigma\xi_{0}^{3}$ may be taken out of the bounce averaging
operator, since $\xi_{0}^{3}$ is an odd function of $\xi_{0}.$ The gradient
vector in the rediced $\left(  p,\xi_{0}\right)  $ space is
\begin{equation}
\nabla_{p,\xi_{0}}=\left(
\begin{array}
[c]{l}%
\nabla_{p}\\
\nabla_{\xi_{0}}%
\end{array}
\right)
\end{equation}
with
\begin{align}
\nabla_{p}  & =\frac{\partial}{\partial p}\\
\nabla_{\xi_{0}}  & =\frac{-\sqrt{1-\xi_{0}^{2}}}{p}\frac{\partial}%
{\partial\xi_{0}}%
\end{align}
